\DTMsavedate{mydate}{2022-06-30}
\maketitle{Modern Control Engineering}{Electronic Engineering}{\DTMusedate{mydate}}

% ----------------------------------------------------- %
% COURSE INFO
\subsection*{Course Information}
\vspace{0ex}%
\noindent\parbox{0.5\textwidth}{%
    \noindent\begin{tabular}{@{}ll}
                 \textsf{Course:}          & 	Modern Control Engineering  \\
                 \textsf{Course Code:}         & 	TEE 5141    \\
                 \textsf{Credit Hours:}    & 	10
    \end{tabular}}
\vspace{2ex}\hrule\vspace{2ex}

% ----------------------------------------------------- %
% INSTRUCTOR INFO
\subsection*{Lecturer Information}
\noindent\parbox{0.5\textwidth}{%
    \noindent\begin{tabular}{@{}ll}
                 \textsf{Name:}         &     Reginald Gonye         \\
%\textsf{Phone:}        &    \phone          \\
\textsf{Email:}        &    reginald.gonye@nust.ac.zw \\
\textsf{Qualifications:}        &  MPhil in Electronic Engineering, (Hon) BEng in Electronic Engineering
    \end{tabular}}
\vspace{2ex}\hrule\vspace{2ex}


% ----------------------------------------------------- %
% COURSE SYNOPSIS
\subsection*{Course Synopsis}
The module looks at State Space Analysis covering State-space methods of analysis and design, Observability and controllability, Pole placement for the optimization response, State observers and pole placement design with state observers, Multi-input, multi-output systems and cross-coupling problems.
It also delves into Digital Control covering Digital time control systems, Modeling of Sampled Processes, Transient response, Steady state response, Stability, Design of Digital Controllers and Root Locus.
\vspace{2ex} \hrule\vspace{2ex}


% ----------------------------------------------------- %
% COURSE DESCRIPTION
\subsection*{Learning Aim}
To instill state space analysis methods for Control Engineering into students.
To instill discrete analysis methods for Control Engineering into students.
\vspace{2ex} \hrule\vspace{2ex}

% ----------------------------------------------------- %
% LEARNING OBJECTIVES
\subsection*{Learning Objectives}
By the end of the course, candidates should be able to:
\begin{outline}
    \1      Understand the concept of states as applied to Control Systems.

    \1  	Be able to use state space analysis in the analysis of Control Systems.

    \1  	Understand the concept of discrete analysis as applied to digital control systems.

    \1  	Be able to analyse digital control systems.

\end{outline}
\vspace{2ex}\hrule\vspace{2ex}

% ----------------------------------------------------- %
% COURSE TOPICS
\subsection*{Topics}
\begin{outline}
    \1 A review of Control Engineering
    \1 Models of Systems
    \1 State Space Representations
    \1 Time Response from State Space Equations
    \1 State Space Realizations
    \1 State Feedback
    \1 Review of analysis of Discrete Signal
    \1 Modelling of Sampled Processes
    \1 Performance e Features of Digital Control Systems
    \1 Design of Digital Controllers
    \1 Root Locus
\end{outline}
\vspace{2ex} \hrule\vspace{2ex}

% ----------------------------------------------------- %
% Students Assenssment
\subsection*{Assessment of Students}
\begin{outline}
    \1 Continuous assessment will consist of tests and assignments (NB: Attendance may also contribute to the coursework mark).
    \1 Final assessment will consist of an examination at the end of the semester. The examination shall comprise of six questions each carrying 25 marks. You will be required to attempt four of the six questions.
\end{outline}
\vspace{2ex}\hrule\vspace{2ex}

% ----------------------------------------------------- %
% Grade Evaluation
\subsection*{Course Evaluation and Grading Policies}
Grades are based on:
Continuous assessment (\qty{25}{\percent}),
Final Exam (\qty{75}{\percent})
\vspace{2ex}\hrule\vspace{2ex}
% ----------------------------------------------------- %
% EQUIPMENT
% https://sites.udel.edu/computing-purchases/personal-specs/
% \subsection*{Essential Equipment}

% \vspace{2ex}\hrule\vspace{2ex}

% Policies and Other information
\subsection*{Policies and Other Information}
% Key Text or some Books
\subsection*{Key Text}

\begin{itemize}
    \item Advanced Control Engineering by R. S. Burns
    \item Introduction to Control Systems by D. K. Anand, R. B. Zmood
    \item Control Engineering by E. A. Parr
\end{itemize}
\vspace{2ex}\hrule
\vspace{2ex}\hrule\vspace{2ex}